\section{Introduction}\label{sec:intro}

Mean--variance optimisation \citep{markowitz1952} is the textbook answer to portfolio
construction, yet in practice its outputs are notoriously unstable: small errors in
estimated means and covariances are amplified by the inversion of an ill-conditioned
covariance matrix into extreme, concentrated and rapidly changing weights
\citep{michaud1989,chopra1993}. Out of sample, optimised portfolios frequently fail to beat
naive equal weighting \citep{demiguel2009}. Two families of remedies have emerged. The first
regularises the inputs, for example by shrinking the covariance matrix
\citep{ledoit2004} or by imposing weight constraints, which \citet{jagannathan2003} show
act as a form of shrinkage. The second abandons optimisation altogether in favour of
heuristic risk allocation, such as equal risk contribution \citep{maillard2010}.

Hierarchical Risk Parity (HRP), introduced by \citet{lopezdeprado2016}, belongs to the
second family and has become one of the most popular allocation methods among
practitioners and in open-source libraries. HRP clusters assets by correlation, orders them
so that similar assets are adjacent, and allocates capital top-down by recursively splitting
each cluster in inverse proportion to its risk. Because it never inverts the covariance
matrix, HRP is robust to estimation noise and can handle singular covariance matrices.

HRP has, however, a practical blind spot: it cannot respect constraints. Real mandates
impose per-asset limits (UCITS concentration rules, single-stock caps), sector or
asset-class caps and floors (investment policy statements), and turnover budgets (to control
trading costs). Recursive bisection has no mechanism for any of these, and in our
multi-asset universe unconstrained HRP places on average 86\% of capital in a single
short-term bond fund, breaching a 15\% single-fund cap and a 40\% asset-class cap at
\emph{every} rebalance. The usual fixes are ad hoc: weights are clipped and renormalised,
which can cycle without ever reaching the feasible set, or projected onto it with a
Euclidean distance that ignores how assets co-move.

This paper proposes \textbf{Constraint-Aware Hierarchical Risk Parity (CA-HRP)}, a
principled way to keep HRP's risk allocation while satisfying any set of linear
constraints exactly. Our contributions are:
\begin{enumerate}
\item \textbf{Capacity-aware recursive bisection.} We define a capacity interval for every
node of the HRP tree and clip each split into it. We prove that the interval is never
empty and that per-asset bounds are met exactly, with no solver
(Proposition~\ref{prop:box}). This yields a solver-free variant, CA-HRP-B.
\item \textbf{Tracking-error projection.} We map the HRP target to the feasible portfolio
with minimum ex-ante tracking error to it. We show the projection is unique, leaves
compliant HRP portfolios unchanged, and redistributes excess weight as a covariance-based
hedge rather than uniformly (Proposition~\ref{prop:te}).
\item \textbf{A transparent evaluation protocol.} The method design was selected on a
validation period only and then frozen; we report a separate held-out test period, all
constraint sets were fixed before any backtest was run, and all code and data are public.
\item \textbf{Evidence on six universes and in simulation.} On 20 multi-asset ETFs
(2008--2026), 43 NIFTY~50 stocks (2011--2026) and 47 US large caps (2006--2026), and on three
external-validation universes added after the method was frozen (Indian ETFs, US sector
funds and global country funds), CA-HRP satisfies every constraint at every rebalance and at
every rebalancing frequency. Where constraints bind, it is the strongest member of the HRP
family; where they rarely bind, it matches HRP. In a Monte Carlo study with a known
covariance it has the highest true Sharpe ratio of all methods.
\item \textbf{A regime-robust variant and two negative results.} We propose a worst-case
tracking-error projection over covariance scenarios with a guarantee over their convex hull
(Proposition~\ref{prop:robust}), and report honestly that it does not improve on CA-HRP,
and that embedding constraints inside the bisection performs worse than projecting
afterwards.
\end{enumerate}

The rest of the paper is organised as follows. Section~\ref{sec:related} reviews related
work, Section~\ref{sec:setting} states the problem, Section~\ref{sec:method} presents the
method and its properties, Section~\ref{sec:design} describes the data and experimental
protocol, Section~\ref{sec:results} reports the main results, Section~\ref{sec:extended}
the extended evaluation, and Section~\ref{sec:conclusion} concludes.

\section{Related Work}\label{sec:related}

\paragraph{Estimation error and robust allocation.} The sensitivity of mean--variance
portfolios to input error is well documented \citep{michaud1989,chopra1993}.
\citet{demiguel2009} find that across many data sets none of fourteen optimised models
consistently beats $1/N$ out of sample. Shrinkage estimators \citep{ledoit2004} reduce the
condition number of the covariance matrix, and \citet{jagannathan2003} show that
no-short-sale constraints behave like covariance shrinkage. Risk-based methods allocate by
risk rather than expected return; the equal-risk-contribution (ERC) portfolio
\citep{maillard2010} equalises marginal risk contributions and can be computed by convex
optimisation \citep{spinu2013}.

\paragraph{Hierarchical allocation.} \citet{lopezdeprado2016} introduced HRP, combining
graph-theoretic clustering with recursive inverse-variance bisection. \citet{raffinot2017}
studied hierarchical clustering-based allocation more broadly, and \citet{raffinot2018}
proposed the Hierarchical Equal Risk Contribution (HERC) portfolio, which splits along the
dendrogram rather than the ordered list. \citet{pfitzinger2019} recognised that HRP lacks
constraint handling and proposed a cluster-based capital allocation that accommodates
weight bounds within the bisection.

\paragraph{Positioning of this work.} Capacity-aware bisection is closest in spirit to
\citet{pfitzinger2019}; our contribution is a capacity definition that also incorporates
group caps, together with a short proof of non-emptiness and exact box feasibility based on
the additivity and sub-additivity of the capacity bounds. More importantly, we show
empirically that enforcing constraints \emph{inside} the bisection has a hidden cost: excess
weight cascades down one branch and can concentrate in volatile assets. Our main method
therefore separates the two concerns: HRP supplies the risk allocation, and a
tracking-error projection \citep{roll1992} enforces arbitrary linear constraints
(group floors, caps and turnover) with minimum risk distortion. To our knowledge this is
the first study to compare constraint-handling strategies for HRP under a common,
pre-registered protocol with a held-out test period.
